\section{Аналіз ситуацій порушення безпеки}


\subsection{Ситуація 1. Передбачувана генерація ключа}

Проблема полягає в тому, що для створення AES-256 ключа використовується поточний час. 
Хоча SHA-256 повертає значення довжиною 256 біт, це не означає, що ключ має 256 біт криптографічної стійкості. 
Якщо зловмисник знає приблизний час створення ключа, він може перебрати можливі значення часу, 
обчислити їхні хеші та знайти потрібний ключ. Таким чином, безпека залежить не від довжини хешу, 
а від кількості можливих початкових значень. Для усунення проблеми необхідно генерувати ключ за допомогою 
криптографічно стійкого генератора випадкових чисел (CSPRNG), який забезпечує достатню непередбачуваність.

\subsection{Ситуація 2. Однаковий seed на різних пристроях}

Якщо однакові пристрої використовують однакове початкове значення для генератора псевдовипадкових чисел, 
вони можуть отримувати однакові послідовності. Якщо ці послідовності використовуються для створення 
криптографічних ключів, різні пристрої можуть мати однакові ключі, що значно знижує безпеку системи. 
У такому випадку компрометація одного пристрою може створити загрозу для інших пристроїв із такими самими ключами. 
Для розв'язання проблеми необхідно використовувати системний CSPRNG, який отримує випадкові дані з надійних джерел операційної системи, 
а не встановлювати однакове фіксоване початкове значення.

\subsection{Ситуація 3. Розголошення правильно згенерованого ключа}

Навіть якщо ключ було створено за допомогою надійного CSPRNG, систему не можна вважати повністю захищеною, 
якщо значення ключа записується до журналу програми. Журнали можуть бути доступними іншим користувачам 
або зберігатися тривалий час, тому зловмисник, який отримає до них доступ, зможе використати ключ для 
розшифрування даних або виконання інших операцій від імені системи. Для усунення проблеми не можна записувати 
секретні ключі до журналів. Якщо ключ уже було розголошено, його потрібно замінити, а доступ до журналів обмежити.


\subsection{Ситуація 4. Недостатньо випадкова генерація RSA-ключа}

Під час створення RSA-ключів джерело випадковості має критично недостатню ентропію, внаслідок чого
різні ключі можуть бути сформовані з пов'язаних або повторно використаних простих чисел. 

Основною проблемою тут є низька якість випадкового матеріалу, який використовується 
для генерації простих множників $p$ та $q$. Якщо у двох незалежних модулів
$N_1 = p_1 \cdot q_1$ та $N_2 = p_2 \cdot q_2$ випадково співпаде хоча б один простий множник
(наприклад, $q_1 = q_2$), зловмисник зможе легко знайти цей спільний дільник за допомогою алгоритму
обчислення найбільшого спільного дільника через формулу $\text{gcd}(N_1, N_2) = q$. Це повністю
компрометує та зламає обидві пари асиметричних ключів, дозволяючи розшифрувати приватний трафік або
підробити цифрові підписи. Способом усунення цієї проблеми є використання виключно апаратно
підкріплених системних криптографічно стійких генераторів (CSPRNG) із високим рівнем ентропії та
попередньою перевіркою згенерованих чисел на унікальність~\cite{SP800-56B, RFC8017}.

\subsection{Ситуація 5. Надмірно тривале використання одного ключа}

Інформаційна система використовує один (\textit{master key}) для захисту даних усіх користувачів протягом
тривалого часу без заміни.

Проблема полягає у відсутності політики ротації ключового матеріалу та порушенні принципів мінімізації
ризиків. Така схема створює загрозу накопичення зашифрованого трафіку для подальшого криптоаналізу та
критично збільшує масштаб наслідків у разі витоку ключа: у разі компрометації єдиного \textit{master key} під
загрозою опиняться абсолютно всі дані користувачів за весь період експлуатації системи. Роль регулярної
ротації ключів полягає у зменшенні обсягу даних, захищених одним ключем, і локалізації наслідків можливої
компрометації, оскільки зловмисник зможе отримати доступ лише до сегмента інформації, зашифрованого
поточним зламаним ключем на конкретному часовому інтервалі~\cite{SP800-57}.
